
\chapter{Introduction}
\begin{refsection}

This chapter served as an introduction to the research topic and highlighted why it was relevant. This chapter provided the study’s background by discussing and identifying the research problems and outlining the research objectives and questions. It also highlighted the study’s significance, scope, and limitations. Furthermore, this chapter included a project dictionary defining key terms and concepts used throughout the study.

\section{Background of the Problem}

In the fast changing digital landscape, the creative industry is seeing an unprecedented increase in demand for automated design solutions. The market for AI-powered design tools is expected to increase from \$5.54 billion in 2024 to \$6.77 billion in 2025, representing a compound annual growth rate (CAGR) of more than 22\% \cite{BusinessResearch2025}. The growth is driven by the requirement for rapid asset development and cross-platform consistency, which aligns with UN Sustainable Development Goal 9, which stresses innovation, and Goal 12, which focuses on minimizing redundant design labor. Maintaining visual brand consistency, such as logos, color palettes, and typography, is no longer an aesthetic choice for businesses but a strategic imperative for establishing consumer trust and loyalty.

Historically, design tools have transitioned from manual, coordinate-based software like Adobe Photoshop to more accessible, template-driven platforms like Canva. This progression toward intelligent design is further exemplified by emerging industry research such as Pomelli by Google Labs. Pomelli utilizes semantic modeling to map high-level aesthetic descriptions directly to component styles, allowing for a more intuitive translation of creative intent \cite {GooglePomelli2024}. However, a significant "Generative Gap" remains. Current industrial tools often treat designs as stochastic or static entities. When a design is exported, it is typically "flattened" into pixel-based formats (JPEG/PNG), which effectively destroys the underlying "Visual DNA"—the mathematical and relational logic of the elements. This leads to Destructive State Management, where any manual refinement made by a user is lost the moment a global theme update is applied.

Furthermore, existing tools frequently suffer from Layout Fragility across Heterogeneous Dimensions. Many industrial tools rely on absolute coordinate-based layouts that lack relational logic. When a user transitions from a vertical format (like an ID) to a horizontal one (like a certificate), the system often fails to determine optimal placement, leading to Responsive Layout Failures (RLFs). These failures include elements protruding off edges or overlapping as layout space becomes scarce, making the output look unprofessional or non-functional \cite{Althomalietal2022,Zerinetal2025}.

Also, while modern Generative AI models can produce high-fidelity imagery, they are often stochastic and computationally "opaque". These models can "hallucinate" artistic elements but lack the Deterministic Vector Structures (SVG) necessary for precise, rule-based orchestration. There is a critical need for a system that treats a design not as a flat file, but as a Stateful Collection of Objects that can maintain visual harmony through mathematical constraints.

Within the context of academic and organizational environments, the need for consistent branding across different common publication materials is a recurring challenge. Without a framework that utilizes Graph Neural Networks (GNN) to understand structural hierarchies and Constraint Satisfaction Problem (CSP) Solvers to maintain layout requirements, producing these assets remains a labor-intensive process prone to human error. Without a shared relational model that encodes an organization's stylistic and structural constraints, even minor brand updates — such as a color refinement or a logo revision — cascade into a full manual redesign cycle across every affected publication material. This gap highlights the pressing need for an automated, constraint-aware approach to brand consistency.

ThemeSync addresses these gaps by proposing a graph-based orchestration framework. By utilizing K-Means Clustering for color extraction and the Scale-Invariant Feature Transform (SIFT) for keypoint analysis, the system can autonomously identify a brand's stylistic markers. By integrating a Bi-directional State Machine, the framework ensures that manual user "nudges" are preserved as "Hard Constraints," allowing for a seamless "Human-in-the-Loop" experience that balances automated efficiency with creative control. By treating every design asset as a node within a relational graph, brand updates reach every format initialized from the master layout.
% \begin{figure}[ht]
%     \centering
% 	\includegraphics[width=0.85\textwidth]{figures/sampleFig1.jpg} 
% 	\caption[Barred spiral galaxy NGC 1300]{Barred spiral galaxy NGC 1300 photographed by Hubble telescope. While the galaxy in the photo is not our sun, it does emit light, much like our sun. Image credit: NASA.}
% 	\label{fig:firstFig}
% \end{figure}

% \begin{table}[ht]
% \centering
% \caption{This is a table}
% \label{tbl:sampleTbl1}
% \resizebox{0.8\textwidth}{!}{%
% \begin{tabular}{llll}
% \hline
% installation & type & capacity (GW) & location \\ \hline
% Longyangxia Dam & photovoltaic & 0.85 & China \\
% Gansu Wind Farm & wind & 6 & China \\
% Sihwa Lake & tidal & 0.254 & South Korea \\ \hline
% \end{tabular}%
% }
% \end{table}

\section{Statement of the Problem}
The rapid transition to multi-platform digital branding demands tools that can seamlessly maintain visual integrity across heterogeneous formats. However, current automated design systems rely heavily on static templates or stochastic generation, creating a significant "Generative Gap" where outputs are rendered as computationally opaque, flat files (e.g., JPEG or PNG). This limitation leads to destructive state management, meaning any manual refinement or "nudge" made by a user is mathematically irreconcilable with the master theme and is subsequently lost during global updates. Furthermore, because these industrial tools utilize absolute coordinate-based layouts rather than relational logic, they suffer from layout fragility; transitioning a design between vastly different aspect ratios frequently results in broken layouts that require extensive, labor-intensive manual correction.

To address these computational bottlenecks, this study developed ThemeSync, a specialized design orchestration framework that transitioned away from immutable files by treating layouts as a stateful collection of objects. The developed method utilized a Constraint Satisfaction Problem (CSP) Solver, specifically the Cassowary Algorithm, to treat manual user edits as "Hard Constraints," triggering real-time recalculations to maintain visual harmony among surrounding elements. Coupled with Graph Neural Networks (GNN) to map structural dependencies and a Bi-directional State Machine to serialize data, the system ensured that local human-in-the-loop refinements and typography updates were synchronized without data loss. 


Ultimately, the purpose of this research was to develop and evaluate this orchestration framework to bridge the gap between automated synthesis and manual refinement. By doing so, the study aimed to answer how Graph-Based DNA Extraction could successfully preserve stylistic markers across varying formats, and to what extent a CSP solver could effectively manage real-time manual edits while maintaining global design rules.

% Enter the statement of the problem here. To cite a study add a bib entry in the references.bib, then use this code \cite{noauthor_biblatex_nodate} to cite the study.

\section{Objectives of the Study}

\subsection{General Objective}

This study aims to design and evaluate ThemeSync, a graph-based design orchestration framework that implements real-time constraint satisfaction and stateful manual refinements across heterogeneous design formats

\subsection{Specific Objectives}

To achieve the general objective, the researchers aim to accomplish the following:

\begin{enumerate}
    \item To design a Relational Logic and Feature Extraction model integrating Graph Neural Networks (GNN), K-Means Clustering, and the SIFT algorithm to autonomously identify and map structural dependencies and stylistic markers from master design files.
    \item To develop a Reactive Layout Synthesis engine utilizing a Bi-directional State Machine and the Cassowary Algorithm to solve Constraint Satisfaction Problems (CSP), enabling the real-time propagation of manual "Hard Constraints" while preserving global layout states.
    \item To evaluate the architectural integrity and professional viability of the generated multi-format layouts by implementing a mathematically calibrated Heuristic Fitness Function to quantitatively rank and optimize the aesthetic validity of parallel design candidates, followed by a quantitative human-expert validation.
\end{enumerate}
   

\section{Significance of the Study}

This study aimed to enhance the traditional design workflow by integrating a graph-based orchestration framework designed to improve the efficiency of producing layouts and to reduce manual labor, and that kept every element within its constraints across formats. This study was intended to benefit the following:

\textbf{Students and Faculty Designers.} Designing school-wide materials such as IDs, certificates, and posters often becomes inefficient due to repetitive resizing and layout inconsistencies. A system for creating "Master Themes" that automatically readjusted themselves to the appropriate size reduced the need for manual formatting of each file, allowing students and faculty to concentrate on subject matter and creative aspects rather than on avoiding technical errors.

\textbf{Camarines Sur Polytechnic Colleges.} The college could benefit from a more efficient process for producing school materials. Automating design tasks would allow the institution to save time and effort on large projects while keeping official documents consistent, accurate, and professional, reducing the operational burden on staff responsible for managing and producing institutional publications.

\textbf{Creative Professionals and Industrial Organization.} Current design tools and production workflows often require manual, repetitive adjustments when resizing assets for various platforms, leading to "Layout Fragility" and high labor costs. This study provided a specialized pipeline for Multi-Format Asset Orchestration that utilized a Constraint Satisfaction Problem (CSP) Solver to automate the scaling of "Visual DNA.”

\textbf{Researchers.} The study contributed to the fields of Intelligent Systems and Design Automation by exploring the intersection of Graph Neural Networks (GNN) and Deterministic State Orchestration. It provided a methodology for Object Serialization and Bi-directional State Machines, offering insights into how design intent—traditionally a subjective human trait—could be mathematically preserved and synchronized within a digital framework.

\textbf{Future Researchers.} The system established a technical foundation that allows future scholars to enhance the automation engine or adapt its algorithms for other digital media. Developers and scholars can utilize the study’s findings on Heuristic Search and Aesthetic Scoring, including the experts' preference for evenly spaced layouts, to build more advanced, user-centric technologies or to investigate real-time collaborative design environments.

\textbf{Computing Field.}  The study applied Artificial Intelligence (AI) and Computer Vision by utilizing K-Means Clustering for color extraction and Heuristic Search Algorithms for layout suggestions. It demonstrated how a Human-in-the-Loop (HITL) approach could combine mathematical precision with human creativity, and it offered an approach to the "Generative Gap" by converting flat, "opaque" files into Stateful Vector Structures, thereby extending the capabilities of smart design environments.

\section{Scope and Limitation}

The primary scope of this research involved the design, development, and technical evaluation of ThemeSync, a graph-based design orchestration framework engineered to automate brand identity synchronization across heterogeneous formats while addressing the challenges of “Generative Gap” and “Layout Fragility” in conventional design tools. Specifically, the study sought to design a Relational Logic and Feature Extraction model that integrated Graph Neural Networks, K-Means Clustering, and the SIFT algorithm to identify and map structural dependencies and stylistic markers from master design files. Furthermore, the scope encompassed the implementation of Dynamic Node Instantiation, enabling the framework to merge extracted ”Visual DNA” with external administrative datasets—like user-defined JSON payloads containing specific names, titles, or signature counts—to generate fully populated, variable-driven documents. It also aimed to implement a Reactive Layout Synthesis engine powered by a Bi-directional State Machine and the Cassowary Algorithm to solve Constraint Satisfaction Problems, ensuring real-time propagation of manual “Hard Constraints” and preservation of global layout states during resizing or format transitions. Lastly, the research intended to evaluate the architectural integrity and aesthetic quality of generated multi-format layouts through heuristic search and a mathematically calibrated fitness function, thereby optimizing structural consistency and visual validity across diverse platforms.

The study was limited to the synchronization of 2D graphical layouts and branding assets; it did not extend to 3D modeling or video-based motion graphics. Furthermore, while the system would provide ”Aesthetic-Aware” suggestions, the final design authority remained with the human user, focusing on a Human-in-the-Loop refinement process rather than autonomous design replacement. Consequently, while the core computational evaluation of ThemeSync relied on rigorous mathematical benchmarking—specifically measuring algorithmic latency, layout detection accuracy via Intersection over Union (IoU), and constraint satisfaction rates—the validation of the final generated layouts did not rely exclusively on theoretical metrics. To ensure that the framework's outputs aligned with real-world professional standards, the study integrated a mixed-methods evaluation strategy. The validated design candidates were then evaluated by a panel of four domain experts (professional layout artists and graphic designers with documented industry experience) using a quantifiable metric. This approach was designed to verify that the framework satisfied both computational constraints and practical industry aesthetic requirements.

\section{Project Dictionary}

 This section provides the definitions of the major terms and concepts that are used throughout this study, serving as a guide in understanding the technologies, methods, and components involved in developing ThemeSync. 

%\begin{itemize}
    \textbf{Affine Transformations}. An affine transformation is the group of transformations that leave the parallelism of straight lines fixed, it can be combined with any transformation consisting of rotation, translation, shearing, and scaling \cite{Torpey&Klein2024}. In this study, Affine Transformations are mathematical functions that are used to transform coordinates and hierarchies of visual images between the master design and new canvas sizes in heterogeneous aspect ratios.

    \textbf{Bi-Directional state machine}. A state machine is a form of automation where the execution of the system operates with a finite amount of states and transitions to another state based on an activity or condition \cite{UiPath2024}. In this study, the calculation algorithm is the one that is used to coordinate the refinements made manually with the automated layouts by sequentializing the refinements into a stable JSON state.

    \textbf{Cassowary algorithm}. Cassowary is an incremental constraint-solving algorithm that efficiently solves systems of linear equalities and inequalities, in which each constraint is either a requirement or a weighted preference; it was developed for user interface layout \cite{Badros2001Cassowary}. In this study, this algorithm is the main layout engine that re-calculates soft constraints when it is injected with manual Hard Constraints to ensure that design elements are not overlapping.

    \textbf{Constraint Satisfaction Problem (CSP).} Constraint Satisfaction Problem (CSP) are core in artificial intelligence and theoretical computer science, a model of variables with domains in finite sets and a list of constraints, which are to be satisfied \cite{Cenetal2025}. In this study, the mathematical model used to control layout density and support recalculation of design nodes in real time to create visual harmony is CSP. This enables the system to dynamically enforce branding constraints across all generated publication materials.

    \textbf{Dynamic Node Instantiation.} Dynamic Node Instantiation is the algorithmic process by which external, user-defined variables are assigned baseline geometric dimensions and converted into functional mathematical nodes. The JSON format acts as a critical standard for lightweight data exchange, and its nested structure is uniquely suited for representing these diverse UI properties \cite{Berdan2025Formats}. In this study, Dynamic Node Instantiation enables seamless communication between the frontend interfaces and the backend Graph Neural Network, allowing dynamic data to be orchestrated alongside the master theme.

    \textbf{Fitness Function.} This is a quantitative figure of merit used to evaluate how effectively a candidate solution satisfies established design composition features and structural requirements \cite{Fitness2025}. In this study, the mathematical equation to rank the layout candidates is the Fitness Function, which is used to determine how well the layouts meet design rules such as alignment to a structural grid and the ratio of whitespace to other elements.

    \textbf{Graph Neural Network (GNN)}. Graph Neural Network (GNN) is a type of neural network that represents graph dependencies by exchanging messages among graph nodes in order to address problems that involve non-Euclidean space \cite{Boetal2023}. In this study, the GNN is the machine learning model that is applied to analyze design hierarchies where elements are viewed as nodes and the relationship between elements as edges to determine how best to position elements.

    \textbf{Hard Constraint.} Hard Constraint is a requirement that must be strictly satisfied in a constraint optimization problem \cite{Williams&Robbins2001}. In this study, any manual user edit, such as moving or resizing an element, is registered as a “Hard Constraint” that the system must respect while adjusting other elements.

    \textbf{Implicit Geometric Hierarchy.} Implicit Geometric Hierarchy is the mathematical principle by which the framework establishes visual priority. Graph Neural Networks are capable of capturing irregular object boundaries and complex spatial relationships with high accuracy by treating distinct areas or patches as individual graph nodes \cite{Han2022VisionGNN, Jiang2024Graph4GUI}. In this study, rather than relying on explicit semantic tags, visual priority is encoded in each element's baseline geometric dimensions and reading order: the Spatial Hierarchy Index sets each element's role and its position in the reading order, and each role is assigned a fixed baseline bounding box, so that the title receives the largest box and the recipient's or attendee's name one of the largest.

    \textbf{K-Means clustering.} A basic activity within data analysis is clustering, which sets out to arrange a sizeable dataset into significant subsets along the similarity measures, with the K-Means algorithm being modest in the sense that the algorithm ascertains the data points and assigns them to the closest cluster centroid \cite{Lertrusdachakul2025Hybrid}. In this study, K-Means Clustering is the algorithm that is utilized in order to automatically identify dominant and accent color palettes in the pixels of a Master Design file.

    \textbf{Mean Opinion Score (MOS).} A standardized, quantitative metric originally formulated to assess the subjective quality and perceptual fidelity of system outputs through structured human evaluation \cite{itut1996p800}. In this study, MOS refers to the structured five-point rating scale rubric utilized by domain experts to quantitatively grade the algorithmically generated multi-format layouts across three formalized evaluation parameters: Visual Hierarchy Preservation, Brand and Typographic Consistency, and Constraint and Whitespace Harmony.

    \textbf{Object Serialization.} Object Serialization is the process of transforming complex data structures into a standardized format to enable real-time synchronization and system architecture persistence \cite{WikipediaSerialization2026}. In this study, Object Serialization refers to the transformation of visual design objects into a structured JSON format to preserve a ”Live State Link” so that these objects can be post-generatively editable.

    \textbf{Semantic Layout Parsing.} Semantic Layout Parsing is the process of identifying and categorizing different parts of a visual layout to understand its underlying structure \cite{BobLd2024}. In this study, it is a pipeline used for PNG inputs that utilizes object detection (YOLO) and OCR to “unflatten” raster images into editable semantic nodes.

    \textbf{SIFT (Scale-Invariant Feature Transform).} SIFT is an algorithm in computer vision used to detect and describe local features in images \cite{Lowe2009}. In this study, SIFT keypoint density is used to rank the decorative motifs extracted from the master design and to build the saliency map against which the legibility of the generated text is scored.

    \textbf{Soft Constraints.} Soft Constraints are preferences or goals within a layout system that are not strictly required for a valid solution but are optimized to achieve the most desirable aesthetic or functional outcome \cite{Jiang2020ORCSolver}. According to Jiang, Stuerzlinger, Zwicker, and Lutteroth (\citeyear{Jiang2020ORCSolver}), these constraints allow for flexible user interface layouts that can adapt to varying screen dimensions while maintaining proportional harmony. In this study, Soft Constraints represent the initial automated placement of design elements which can be overriden or re-calculated when a user introduces manual Hard Constraints.    

    \textbf{Visual DNA.} An algorithm to encode images as distributions of neuron activations (DNAs) to form an expressive, high-dimensional blueprint enabling visual attributes in one dataset to be compared and preserved in another dataset \cite{Matzetal2023}. In this study, Visual DNA refers to a hybrid data structure combining autonomously extracted algorithmic features (color centroids and relational logic) with user-defined administrative constraints (specific font families, weights, and baseline sizes).

    % \textbf{Pomelli.} Pomelli is an experimental design system that utilizes semantic modeling to map natural language aesthetic descriptions directly to UI component styles. As established by Google Labs \cite{GooglePomelli2024}, this framework enables the intuitive translation of creative intent into visual assets. In this study, Pomelli serves as a contextual benchmark for semantic style transfer, which the proposed framework builds upon by adding graph-based structural permanence and real-time constraint satisfaction.
%\end{itemize}

%=======================================================%
%%%%% Do not delete this part %%%%%%
\clearpage

\begin{center}
\textbf{Notes}
\end{center}
\printbibliography[heading=none]
\end{refsection}